
%%%%%%%%%%%%%%%%%%%%%%%%%%%%%%%%%%%%%%%%
%           main body
%%%%%%%%%%%%%%%%%%%%%%%%%%%%%%%%%%%%%%%%

%-----------------------------------
%	TITLE PAGE
%-----------------------------------

\title[数学分析]{\texorpdfstring{\LARGE \bfseries 第九章\quad 数项级数}{第九章 数项级数}} % The short title appears at the bottom of every slide, the full title is only on the title page
\author[任意项级数] % Your institution as it will appear on the bottom of every slide, maybe shorthand to save space
{\Large \bfseries 任意项级数}
% \author{Singular Value Decomposition} % Your name
% \date{\today} % Date, can be changed to a custom date
\date{}

%------------------------------------------------

\begin{document}

%------------------------------------------------

\begin{frame}

\titlepage % Print the title page as the first slide

\end{frame}

%------------------------------------------------

% \begin{frame}
% \frametitle{内容提要} % Table of contents slide, comment this block out to remove it
% \tableofcontents % Throughout your presentation, if you choose to use \section{} and \subsection{} commands, these will automatically be printed on this slide as an overview of your presentation
% \end{frame}

%-----------------------------------
%	PRESENTATION SLIDES
%-----------------------------------

%------------------------------------------------

\begin{frame}{从正项级数到任意项级数}

\begin{columns}
\begin{column}{0.55\textwidth}
一个级数, 如果只有有限个负项或有限个正项, 都可以用正项级数的~{\color{red}各种判别法}来判断它的敛散性.

\vspace{0.6em}

\only<1-6>{\phantom{如果一个级数即有无限个正项, 又有无限个负项, 那么正项级数的各种判别法不再适用！}}
\only<7->{如果一个级数即有无限个正项, 又有无限个负项, 那么正项级数的各种判别法不再适用！}
\end{column}

\begin{column}{0.42\textwidth}
\only<2->{
\begin{empheq}[box={\fancymath[colback=cyan!30, sharp corners]}]{align*}
  & \only<2->{\text{\color{red!80}\ding{43} 比较判别法}} \\
  & \only<3->{\text{\color{red!80}\ding{43} Cauchy判别法}} \\
  & \only<4->{\text{\color{red!80}\ding{43} d'Alembert判别法}} \\
  & \only<5->{\text{\color{red!80}\ding{43} Raabe判别法}} \\
  & \only<6->{\text{\color{red!80}\ding{43} 积分判别法}} \\
  & \only<6->{\text{\color{red!80}\ding{43} $\cdots\cdots$}}
\end{empheq}
}
\end{column}
\end{columns}

\vspace{1em}

\only<8->{
\begin{center}
我们只能从最基本的~{\color{red}Cauchy收敛原理}出发, 进行研究.
\end{center}
}

\end{frame}

%------------------------------------------------

\section[Cauchy\\收敛原理]{Cauchy收敛原理}

%------------------------------------------------

\begin{frame}{级数的Cauchy收敛原理}

\begin{thm}[Cauchy收敛原理]
级数$\sum\limits_{n=1}^\infty x_n$收敛的充分必要条件是：对任意给定的$\varepsilon > 0,$ 存在正整数$N,$ 使得
\[\left\lvert x_{n+1} + x_{n+2} + \cdots + x_m \right\rvert = \left\lvert \sum\limits_{k=n+1}^m x_k \right\rvert < \varepsilon\]
对一切$m > n > N$成立.
\end{thm}

\pause
\vspace{0.6em}

\begin{attnblock}
\begin{itemize}
\item[\ding{43}] Cauchy收敛原理对一般的度量空间都是成立的
\pause
\item[\ding{43}] 取$m = n + 1,$ 我们即得到级数收敛的必要条件$\lim\limits_{n\to\infty} x_n = 0.$
\end{itemize}
\end{attnblock}

\end{frame}

%------------------------------------------------

\section[Leibniz\\级数]{Leibniz级数}

%------------------------------------------------

\begin{frame}{Leibniz级数}

\begin{Def}[交错级数]
如果级数$\sum\limits_{n=1}^\infty x_n = \sum\limits_{n=1}^\infty (-1)^{n+1} u_n ~ (u_n > 0),$ 则称其为\textbf{交错级数}.
\end{Def}

\pause

\begin{Def}[Leibniz级数]
如果级数$\sum\limits_{n=1}^\infty x_n = \sum\limits_{n=1}^\infty (-1)^{n+1} u_n ~ (u_n > 0)$为交错级数, 并且数列$\{u_n\}$单调减少且收敛于$0,$ 则称此级数为\textbf{Leibniz级数}.
\end{Def}

\pause

\begin{block}{Leibniz级数的一些例子}
\vspace{-1.5em}
\[\sum\limits_{n=1}^\infty \dfrac{(-1)^{n+1}}{n^s} (s > 0), ~~~ \sum\limits_{n=2}^\infty (-1)^{n} \dfrac{\ln{n}}{n}, ~~~ \sum\limits_{n=1}^\infty (-1)^{n+1} \dfrac{n^2}{n^3 + 1}, ~~ \cdots\]
\end{block}

\end{frame}

%------------------------------------------------

\begin{frame}{Leibniz级数}

有一些Leibniz级数, 则不那么明显能一眼看出：

\pause
\vspace{0.5em}

考虑级数$\sum\limits_{n=1}^\infty \sin \left( \sqrt{n^2+1} \pi \right),$ 有
\begin{align*}
\sin \left( \sqrt{n^2+1} \pi \right) & = {\color{red} (-1)^n} \sin \left( \sqrt{n^2+1} \pi ~ {\color{red} - n\pi} \right) \\
& = (-1)^n \sin \left( \sqrt{n^2+1} - n \right) \pi \\
& = (-1)^n ~ \underbrace{\color{red} \sin \dfrac{\pi}{\sqrt{n^2+1} + n}}_{u_n}.
\end{align*}

\pause

容易看出$\{u_n = \sin \dfrac{\pi}{\sqrt{n^2+1} + n}\}$是单调减小数列, 且
\[\lim\limits_{n\to\infty} \sin \dfrac{\pi}{\sqrt{n^2+1} + n} = 0,\]
可知$\sum\limits_{n=1}^\infty \sin \left( \sqrt{n^2+1} \pi \right)$是Leibniz级数.

\end{frame}

%------------------------------------------------

\begin{frame}{Leibniz判别法}

\begin{thm}[Leibniz判别法]
\begin{center}
Leibniz级数必定收敛.
\end{center}
\end{thm}

\pause
\vspace{0.5em}

\textbf{证明：} 首先, 记$p = m - n,$ 有
\[\left\lvert x_{n+1} + x_{n+2} + \cdots + x_{n+p} \right\rvert = \left\lvert u_{n+1} - u_{n+2} + \cdots + (-1)^{p+1} u_{n+p} \right\rvert.\]
\pause
当$p$是奇数时, 有
\begin{multline*}
u_{n+1} - u_{n+2} + \cdots + (-1)^{p+1} u_{n+p} \\
= \begin{cases}
(u_{n+1} - u_{n+2}) + \cdots + (u_{n+p-2} - u_{n+p-1}) + u_{n+p} ~ {\color{red} \geqslant u_{n+p}}, \\
u_{n+1} - (u_{n+2} - u_{n+3}) - \cdots - (u_{n+p-1} - u_{n+p}) ~ {\color{red} \leqslant u_{n+1}}
\end{cases}
\end{multline*}

\end{frame}

%------------------------------------------------

\begin{frame}{Leibniz判别法的证明}

当$p$是偶数时, 有
\begin{multline*}
u_{n+1} - u_{n+2} + \cdots + (-1)^{p+1} u_{n+p} \\
= \begin{cases}
(u_{n+1} - u_{n+2}) + \cdots + (u_{n+p-2} - u_{n+p-1}) ~ {\color{red} \geqslant 0}, \\
u_{n+1} - (u_{n+2} - u_{n+3}) - \cdots - u_{n+p} ~ {\color{red} \leqslant u_{n+1}}.
\end{cases}
\end{multline*}

\pause

总之, 成立
\begin{align*}
\left\lvert x_{n+1} + x_{n+2} + \cdots + x_{n+p} \right\rvert & = \left\lvert u_{n+1} - u_{n+2} + \cdots + (-1)^{p+1} u_{n+p} \right\rvert \\
& \leqslant u_{n+1}.
\end{align*}

由于$\lim\limits_{n\to\infty} u_n = 0,$ 所以对于任意给定的$\varepsilon > 0,$ 存在正整数$N,$ 使得对一切$n > N,$ 成立
\only<2-2>{
\[\varepsilon > u_{n+1},\]
}
\only<3->{
\[\varepsilon > u_{n+1} ~ {\color{red} \geqslant \left\lvert x_{n+1} + x_{n+2} + \cdots + x_{n+p} \right\rvert},\]
}
\only<4->{
由Cauchy收敛原理知Leibniz级数必收敛.
}

\end{frame}

%------------------------------------------------

\section[A-D判别法]{Abel判别法与Dirichlet判别法}

%------------------------------------------------

\begin{frame}{回忆 | {\smaller[2] 反常积分的A-D判别法}}

我们先回忆一下一般函数的反常积分敛散性的Abel判别法与Dirichlet判别法 (A-D判别法)：

\pause

考虑$\int_a^{+\infty} f(x)g(x) \mathrm{d} x,$ 有
\begin{itemize}
\item[(1)] Abel判别法：
\[\int_a^{+\infty} f(x) \mathrm{d} x ~ \text{收敛}, ~ g(x) ~ \text{单调有界} \Longrightarrow \int_a^{+\infty} f(x)g(x) \mathrm{d} x ~ \text{收敛};\]
\item[(2)] Dirichlet判别法：
\begin{multline*}
\int_a^{A} f(x) \mathrm{d} x ~ \text{关于$A$有界}, ~ g(x) ~ \text{单调}, ~ \lim\limits_{x\to\infty} g(x) = 0 \\
\Longrightarrow \int_a^{+\infty} f(x)g(x) \mathrm{d} x ~ \text{收敛}
\end{multline*}
\end{itemize}

\pause

接下来要考虑的一般级数, 也假设它具有$\sum\limits_{n=1}^{\infty} a_n b_n$的形式, 我们有类似的结论, 证明方法也是类似的.

\end{frame}

%------------------------------------------------

\begin{frame}{级数的A-D判别法}

\begin{thm}[级数的A-D判别法]
若下列条件之一满足, 则级数$\sum\limits_{n=1}^{\infty} a_n b_n$收敛：
\begin{itemize}
\item[(1)] Abel判别法：
\[\sum\limits_{n=1}^{\infty} b_n ~ \text{收敛}, ~~~ a_n ~ \text{单调有界}\]
\item[(2)] Dirichlet判别法：
\[\sum\limits_{n=1}^{N} b_n ~ \text{关于$N$有界}, ~~~ a_n ~ \text{单调趋于$0$}\]
\end{itemize}
\end{thm}

\end{frame}

%------------------------------------------------

\begin{frame}{A-D判别法的分析与证明}

对于一般的情况, 目前能用的工具只有Cauchy收敛原理, 所以要考察
\[\left\lvert \sum\limits_{k=n+1}^{m} a_k b_k \right\rvert.\]

\pause

(1). Abel判别法：条件为$\sum\limits_{n=1}^{\infty} b_n$收敛, $a_n$单调有界.
\vspace{-1em}
\begin{align*}
\sum\limits_{n=1}^{\infty} b_n ~ \text{收敛} ~~ & ~ {\color{red}\Longrightarrow} ~~
\only<2>{\phantom{\forall~\varepsilon > 0, \exists N, \text{ 使得 } \forall~m > n > N, \left\lvert \sum\limits_{k=n+1}^{m} b_k \right\rvert < \varepsilon.}}
\only<3->{\forall~\varepsilon > 0, \exists N, \text{ 使得 } \forall~m > n > N, \left\lvert \sum\limits_{k=n+1}^{m} b_k \right\rvert < \varepsilon.} \\
\only<4->{a_n ~ \text{单调有界} ~~ & ~ {\color{red}\Longrightarrow}} ~~ \only<5->{\exists M > 0, \text{ 使得 } \lvert a_n \rvert < M \text{ 对所有$n$成立}}
\end{align*}

\vspace{-1em}

\only<6->{
\begin{empheq}[box={\fancymath[colback=cyan!30, drop lifted shadow, sharp corners]}]{equation*}
  \text{\bfseries\color{red!80}希望： $\sum\limits_{k=n+1}^{m} a_k b_k \rightsquigarrow \left( \sum\limits_{k=n+1}^{m} b_k \right) \cdot \big( a_{n+1}, \cdots, a_m \text{ 相关的项} \big)$}
\end{empheq}
}

\end{frame}

%------------------------------------------------

\begin{frame}{A-D判别法的分析与证明}

事实上, 如果记$B_t = \sum\limits_{k=n+1}^{t} b_k, t > n,$ 并约定~${\color{red}B_n = 0},$ 那么所有的$B_t$都满足$\lvert B_t \rvert < \varepsilon,$ 并且有
{\smaller[1]
\begin{align*}
\left\lvert \sum\limits_{k=n+1}^{m} a_k b_k \right\rvert & = \left\lvert \sum\limits_{k=n+1}^{m} a_k \left( B_{k} - B_{k-1} \right) \right\rvert = \left\lvert \sum\limits_{k=n+1}^{m} a_k B_{k} - \sum\limits_{k = ~ {\color{red}n+1}}^{{\color{red}m}} a_{{\color{red}k}} B_{{\color{red}k-1}} \right\rvert \\
& = \left\lvert \sum\limits_{k=n+1}^{m} a_k B_{k} - \sum\limits_{k = ~ {\color{red}n}}^{{\color{red}m-1}} a_{{\color{red}k+1}} B_{{\color{red}k}} \right\rvert \\
& = \left\lvert a_mB_m - \sum\limits_{k=n+1}^{m-1} (a_{k+1} - a_k) B_{k} + a_{n+1}{\color{red}B_n} \right\rvert \\
& \leqslant \left\lvert a_m \right\rvert \varepsilon + \sum\limits_{k=n+1}^{m-1} \left\lvert a_{k+1} - a_k \right\rvert \varepsilon
\end{align*}
}

\end{frame}

%------------------------------------------------

\begin{frame}{A-D判别法的分析与证明}

由于$a_n$单调, 所以$\sum\limits_{k=n+1}^{m-1} \left\lvert a_{k+1} - a_k \right\rvert = \left\lvert a_{n+1} - a_{m} \right\rvert,$ 于是
\begin{align*}
\left\lvert \sum\limits_{k=n+1}^{m} a_k b_k \right\rvert & \leqslant \left\lvert a_m \right\rvert \varepsilon + {\color{red} \sum\limits_{k=n+1}^{m-1} \left\lvert a_{k+1} - a_k \right\rvert} \varepsilon \\
& = \left( \left\lvert a_m \right\rvert + {\color{red} \left\lvert a_{n+1} - a_{m} \right\rvert} \right) \varepsilon \\
& < 3M\varepsilon
\end{align*}
\pause
有Cauchy收敛原理, 阿贝判别法得证.

\end{frame}

%------------------------------------------------

\begin{frame}{A-D判别法的分析与证明}

(2). Dirichlet判别法：条件为$\sum\limits_{n=1}^{N} b_n$关于$N$有界, $a_n$单调趋于$0.$
\begin{align*}
a_n ~ \text{单调趋于$0$} ~~ & {\color{red}\Longrightarrow} ~~
\only<1>{\phantom{\forall~\varepsilon > 0, \exists N, \text{ 使得 } \forall~n > N, \left\lvert a_n \right\rvert < \varepsilon.}}
\only<2->{\forall~\varepsilon > 0, \exists N, \text{ 使得 } \forall~n > N, \left\lvert a_n \right\rvert < \varepsilon.} \\
\only<3->{\sum\limits_{n=1}^{N} b_n ~ \text{关于$N$有界} ~~ & {\color{red}\Longrightarrow}} ~~ \only<4->{\exists M > 0, \text{ 使得 } \forall~N \text{ 有 } \left\lvert \sum\limits_{n=1}^{N} b_n \right\rvert < ~ {\color{red}M}}
\end{align*}

\only<5->{
利用证明阿贝判别法类似的方法, 有
\begin{align*}
\left\lvert \sum\limits_{k=n+1}^{m} a_k b_k \right\rvert & = \left\lvert a_m ~ {\color{red}B_m} - \sum\limits_{k=n+1}^{m-1} (a_{k+1} - a_k) ~ {\color{red}B_{k}} \right\rvert \\
& \leqslant \left\lvert a_m \right\rvert \cdot ~ {\color{red}2M} + \sum\limits_{k=n+1}^{m-1} \left\lvert a_{k+1} - a_k \right\rvert \cdot ~ {\color{red}2M} \\
& = \left( \left\lvert a_m \right\rvert + \left\lvert a_{n+1} - a_{m} \right\rvert \right) \cdot 2M \\
& < 6M\varepsilon
\end{align*}
}

\end{frame}

%------------------------------------------------

\begin{frame}{A-D判别法的分析与证明}

反常积分敛散性的A-D判别法 (实际是证明过程中用到的积分第二中值定理) 与级数敛散性的A-D判别法的证明实际上都用到了如下的和式转换方法

\begin{lem}[Abel变换 (分部求和公式)]
设$\{a_n\}, \{b_n\}$是数列, 记$B_n = \sum\limits_{k=1}^n b_k,$ $n = 1, 2, \dots,$ 那么
\[\sum\limits_{k=1}^n a_kb_k = a_n B_n - \sum\limits_{k=1}^{n-1} (a_{k+1} - a_k) B_k.\]
\end{lem}

\end{frame}

%------------------------------------------------

\begin{frame}{A-D判别法的分析与证明}

相应地, 我们对$\sum\limits_{k=1}^n a_kb_k$有如下的估计

\begin{lem}[Abel引理]
设$\{a_n\}$为单调数列, $\{B_n = \sum\limits_{k=1}^n b_k\}$为有界数列, $M_0$为它的一个界, 那么
\[\left\lvert \sum\limits_{k=1}^n a_kb_k \right\rvert \leqslant M_0 \cdot \left( \left\lvert a_n \right\rvert + \left\lvert a_n - a_1 \right\rvert \right).\]
\pause
特别地, 如果进一步有$m \leqslant B_n \leqslant M,$ 并且$\{a_n\}$非负不增, 那么
\[m a_1 \leqslant \sum\limits_{k=1}^n a_kb_k \leqslant M a_1\]
\end{lem}

\end{frame}

%------------------------------------------------

\begin{frame}{A-D 判别法的几个注记}

1. 回过头来看 Leibniz 判别法：
\begin{center}
Leibniz 级数都收敛
\end{center}
\pause
\begin{empheq}[box={\fancymath[colback=cyan!30, drop lifted shadow, sharp corners]}]{equation*}
\text{\bfseries\color{red!80}Leibniz级数： $\sum\limits_{n=1}^{\infty} (-1)^{n+1} u_n, ~~ 0 < u_n \text{ 单调减小, 且 } u_n \to 0.$}
\end{empheq}

\pause
\vspace{0.6em}

若令 $b_n = (-1)^{n+1}, a_n = u_n,$ 那么它们满足 Dirichlet 判别法的条件
\begin{empheq}[box={\fancymath[colback=cyan!30, drop lifted shadow, sharp corners]}]{equation*}
\color{red!80} \sum\limits_{n=1}^{N} b_n ~ \text{关于 $N$ 有界}, ~~~ a_n ~ \text{单调趋于 $0$}
\end{empheq}

\pause

也就是说, Leibniz 判别法可以由 Dirichlet 判别法导出, 可视为是 Dirichlet 判别法的一个特例.

\end{frame}

%------------------------------------------------

\begin{frame}{A-D 判别法的几个注记}

2. Abel 判别法可以由 Dirichlet 判别法导出：

\pause
\vspace{0.8em}

事实上, 若 Abel 判别法的条件成立: $\sum\limits_{n=1}^{\infty} b_n$ 收敛, $a_n$ 单调有界. \par
由数列的单调有界定理, $a_n$ 收敛于某个实数 $a = \lim\limits_{n\to\infty} a_n.$ 那么 $c_n : = a_n - a$ 就构成了一个单调趋于 $0$ 的数列.

\pause
\vspace{0.8em}

另一方面, 由于 $\sum\limits_{n=1}^{\infty} b_n$ 收敛, 那么 $\sum\limits_{n=1}^{N} b_n$ 显然关于 $N$ 有界. 于是, 根据 Dirichlet 判别法, $\sum\limits_{n=1}^{\infty} c_nb_n$ 收敛, 从而知
\[\sum\limits_{n=1}^{\infty} a_nb_n = \sum\limits_{n=1}^{\infty} c_nb_n + a \sum\limits_{n=1}^{\infty} b_n\]
也是收敛的.

\end{frame}

%------------------------------------------------

\begin{frame}{A-D 判别法的应用}

\begin{eg}
设数列 $\{a_n\}$ 单调趋于 $0,$ 则 $\forall~x \in \mathbb{R},$ 级数 $\sum\limits_{n=1}^\infty a_n \sin nx$ 收敛.
\end{eg}

\pause
\vspace{0.4em}

\textbf{证明:} $\{a_n\}$ 单调趋于 $0,$ 这一条件是 Dirichlet 判别法的一半; 自然希望证明另一半条件, 即 $\sum\limits_{n=1}^N \sin nx$ 关于 $N$ 有界.

\pause

\textbf{方法一:} 利用 $\sin nx = \dfrac{e^{inx} - e^{-inx}}{2i},$ 可以将 $\sum\limits_{n=1}^N \sin nx$ 表示为两个等比数列求和 ($x \neq 2k\pi$)：
\begin{align*}
\sum\limits_{n=1}^N \sin nx & = \dfrac{1}{2i} \left( \sum\limits_{n=1}^N \left(e^{ix}\right)^n - \sum\limits_{n=1}^N \left(e^{-ix}\right)^n \right) \\
& = \dfrac{1}{2\sin\frac{x}{2}} \left( \cos\dfrac{x}{2} - \cos \left(N + \dfrac{1}{2}\right) x \right)
\end{align*}

\end{frame}

%------------------------------------------------

\begin{frame}{A-D 判别法的应用}

\textbf{方法二:} 利用三角函数的和差化积、积化和差公式：
{
% \smaller[1]
\begin{align*}
2\sum\limits_{n=1}^N \sin nx & = \sin x + \left( \sin x + \sin 2x \right) + \cdots + \sin Nx \\
& = \sin x + ~ {\color{red} \sin \dfrac{3}{2}x} \cos \dfrac{1}{2}x + \cdots + ~ {\color{red} \sin \dfrac{2N-1}{2}x} \cos \dfrac{1}{2}x + \sin Nx
\end{align*}
}
\pause
两边同时乘以 ${\color{red} \sin \dfrac{1}{2}x},$ 再利用
\[\sin \dfrac{1}{2}x \sin \sin \dfrac{2k-1}{2}x = \cos (k-1)x - \cos kx,\]
相邻项交错相消, 有
{
% \smaller[1]
\begin{align*}
2\sin \dfrac{1}{2}x\sum\limits_{n=1}^N \sin nx
= & \sin \dfrac{1}{2}x \sin x + \cos \dfrac{1}{2}x \cos x - \cos \dfrac{1}{2}x \cos Nx \\
& + \sin \dfrac{1}{2}x \sin Nx
= \cos\dfrac{x}{2} - \cos \left(N + \dfrac{1}{2}\right) x
\end{align*}
}

\end{frame}

%------------------------------------------------

\begin{frame}{A-D 判别法的应用}

总之, 当 $x \neq 2k\pi$ 时, 有
\[\sum\limits_{n=1}^N \sin nx = \dfrac{1}{2\sin\frac{x}{2}} \left( \cos\dfrac{x}{2} - \cos \left(N + \dfrac{1}{2}\right) x \right),\]
对于固定的 $x,$ 上式关于 $N$ 显然是有界的, $\dfrac{1}{\left\lvert \sin\frac{x}{2} \right\rvert}$ 即为它的一个界. 于是, 根据 Dirichlet 判别法, $\sum\limits_{n=1}^\infty a_n \sin nx$ 收敛.

\pause
\vspace{0.6em}

当 $x = 2k\pi$ 时, $\sum\limits_{n=1}^\infty a_n \sin nx$ 的每一项都是 $0,$ 显然也是收敛的.

\end{frame}

%------------------------------------------------

\section{总结}

%------------------------------------------------

\begin{frame}{总结}

{\larger[2] \color{red} \ding{43}} 级数的Cauchy收敛原理 (充要条件)：
\[\left\lvert x_{n+1} + x_{n+2} + \cdots + x_m \right\rvert = \left\lvert \sum\limits_{k=n+1}^m x_k \right\rvert < \varepsilon\]

\pause

{\larger[2] \color{red} \ding{43}} Leibniz判别法：Leibniz级数都收敛.
\[\text{Leibniz级数:}~\sum\limits_{n=1}^{\infty} (-1)^{n+1} u_n, ~~ 0 < u_n \text{ 单调减小, 且 } u_n \to 0.\]

\pause

{\larger[2] \color{red} \ding{43}} A-D 判别法 ($\sum\limits_{n=1}^{\infty} a_n b_n$收敛的充分条件)：
\begin{align*}
\text{Abel判别法:} ~ & \sum\limits_{n=1}^{\infty} b_n ~ \text{收敛}, ~ a_n ~ \text{单调有界} \\
\text{Dirichlet判别法:} ~ & \sum\limits_{n=1}^{N} b_n ~ \text{关于 $N$ 有界}, ~ a_n ~ \text{单调趋于 $0$}
\end{align*}

\end{frame}

%------------------------------------------------

%------------------------------------------------
% % closing page
% \begin{frame}
% \HUGE{\centerline{\bfseries {\color{gray} The} {\color{zkblue} End}}}

% \pause

% \vspace{1.2em}

% \Large{\centerline{\bfseries {\color{gray}Thank} \quad {\color{zkblue} You}}}

% \end{frame}

%----------------------------------------------------------------------------------------

\end{document}
